% Generated draft: requires mathematical and source review.
\documentclass[11pt]{article}
\usepackage[margin=25mm]{geometry}
\usepackage{fontspec}
\setmainfont{Latin Modern Roman}
\usepackage{amsmath,amssymb,mathtools,unicode-math}
\setmathfont{Latin Modern Math}
\usepackage{enumitem,needspace,xurl,hyperref}
\setlength{\parindent}{0pt}
\setlength{\parskip}{5pt}
\setlength{\emergencystretch}{3em}
\newcommand{\N}{\mathbb{N}}
\newcommand{\Z}{\mathbb{Z}}
\newcommand{\Q}{\mathbb{Q}}
\newcommand{\R}{\mathbb{R}}
\newcommand{\C}{\mathbb{C}}
\newcommand{\F}{\mathbb{F}}
\newcommand{\A}{\mathbb{A}}
\newcommand{\PP}{\mathbb{P}}
\newcommand{\E}{\mathbb{E}}
\newcommand{\eps}{\varepsilon}
\newcommand{\dd}{\,\mathrm{d}}
\begin{document}
\section*{UnsolvedMath ID: 30001395}
\textbf{Source problem code:} catalogue entry
\subsection{Definitions and mathematical statement}
The supplied record contains two questions. Its cleaned statement retains only the first; both are preserved here.

For the first question, identify \(\R^2\) with \(\C\). A nonconstant sense-preserving map \(p\colon\R^2\to\R^2\) is \(K\)-quasiregular if \(p\in W^{1,2}_{\mathrm{loc}}\) and
\[
\|Dp(x)\|_{\mathrm{op}}^2\leq K J_p(x)
\quad\text{for almost every }x\in\R^2,
\]
where \(K\geq1\), \(Dp\) is the weak derivative, and \(J_p=\det Dp\). It is of polynomial type if
\[
|p(x)|\longrightarrow\infty\qquad\text{as }|x|\longrightarrow\infty.
\]
Such a map extends continuously to \(\widehat{\R^2}=\R^2\cup\{\infty\}\) by fixing infinity; write \(d=\deg p\) for its positive topological degree. Define \(p^0=\operatorname{id}\) and \(p^{n+1}=p\circ p^n\).

The record does not define the Fatou set for these quasiregular maps. To display the question without silently choosing a convention, let \(\mathcal F(p)\) denote the intended Fatou set. Under a convention in which forward images of its connected components lie in connected components of \(\mathcal F(p)\), a component \(U\) is wandering if there are pairwise distinct connected components \(U_n\) of \(\mathcal F(p)\), for \(n\geq0\), such that
\[
U_0=U,\qquad p^n(U)\subseteq U_n\quad(n\geq0).
\]
\textbf{Question 1.} Can there exist a \(K\)-quasiregular map \(p\colon\R^2\to\R^2\) of polynomial type, with \(K<d=\deg p\), having a wandering Fatou component? The intended quasiregular Fatou-set convention must be confirmed before this becomes an unambiguous statement; in particular, no uniform bound on the dilatations of all iterates is supplied or assumed.

For the additional question in the original statement, let \(X\) and \(Y\) be compact connected Riemann surfaces, with \(X\) hyperbolic, equivalently of genus at least two, and let \(W\subseteq Y\) be a domain, meaning a nonempty connected open subset. A simple island of a holomorphic map \(f\colon W\to X\) over a Jordan domain \(V\subset X\) is a domain \(D\) with compact closure contained in \(W\) such that \(f|_D\colon D\to V\) is a conformal isomorphism.

One usual interpretation of the Ahlfors islands property is the following: \(f\) is nonconstant and holomorphic, and there exists an integer \(k\geq1\) such that, for every collection of Jordan domains \(V_1,\ldots,V_k\subset X\) with pairwise disjoint closures, every connected open set \(O\subseteq Y\) meeting \(\partial_Y W\), and every connected component \(O_0\) of \(O\cap W\), there are \(j\in\{1,\ldots,k\}\) and a domain \(D\) with compact closure contained in \(O_0\) for which \(f|_D\colon D\to V_j\) is a conformal isomorphism. The supplied record does not give its definition of an Ahlfors islands map, so this interpretation requires confirmation against the source.

\textbf{Question 2.} For which triples \((X,Y,W)\) satisfying the preceding surface hypotheses does an Ahlfors islands map \(f\colon W\to X\) exist? How does existence depend on the choice of the hyperbolic target \(X\)? The record does not specify whether \(W\) must be proper or what precise comparison of targets is intended. If \(W=Y\), the displayed boundary condition is vacuous, so this boundary case must be resolved explicitly.
\subsection{Short English statement}
The main question asks whether a polynomial-type quasiregular plane map whose degree exceeds its distortion bound can have a wandering Fatou component. The original record also asks when Ahlfors islands maps to compact hyperbolic Riemann surfaces exist and how existence depends on the target. The Fatou-set convention, the precise islands-map definition, and the allowed boundary cases need confirmation.
\subsection{Sources}
\begin{itemize}
\item Local supplied registry record: \texttt{data/unsolvedmath/problems.json}, permanent UnsolvedMath ID \(30001395\), source problem code \texttt{catalogue entry}, title \emph{Wandering Domains of Polynomial-Type Quasiregular Maps}; collection: \emph{Oberwolfach Reports Open Problems}. No individual proposer is supplied. The proposed year is recorded as \(2009\).
\item Supplied report citation: \emph{Mini-Workshop: The Escaping Set in Transcendental Dynamics} (2010), \emph{Oberwolfach Reports}, pp.\ 2927--2964, DOI \texttt{10.4171/owr/2009/54}. \url{https://doi.org/10.4171/owr/2009/54}. The background describes these reports as documenting mathematical workshops and their open-problem discussions.
\item Supplied literature-source label: ``Rempe Rippon exotic Ahlfors islands.'' \url{https://arxiv.org/abs/1008.1724}. The record supplies no fuller bibliographic citation for this item.
\item Supplied literature-source label: ``original OWR report.'' \url{https://doi.org/10.4171/owr/2009/54}.
\item Literature assessment dated exactly \texttt{2026-08-21}: the supplied curation labels the record \texttt{partially\_solved} and states that Ahlfors-islands existence and wandering-domain questions have constructions for the sphere, torus, and some hyperbolic targets, while the general existence criterion and dependence on the target remain unresolved. This is supplied dated curation, not an independently verified assessment. The sphere and torus fall outside the hyperbolic-target hypothesis of Question 2; the assessment does not explicitly establish the status of Question 1.
\item The registry labels the cleaned statement \texttt{corrected\_verified} and reports that its content was checked against the official Oberwolfach report, with notation, context, and question form normalized. This verification claim is supplied metadata and has not been independently checked here.
\end{itemize}
\end{document}
